% 【本文件写什么】impl 实现层：kernel
% - 定位：内核内建伪 FS / rootfs 逻辑，devfs、procfs、rootfs。
% - crate 路径：os/components/wateros-fs/fs-rootfs/rootfs-impl/impl-kernel/src/
% - 如何实现对应 api-v0 trait：关键类型、状态机、数据结构
% - 算法或硬件交互流程，可用伪代码/流程图
% - 本 impl 启用的 Cargo [features] 与 arch feature，impl-riscv64 等
% - 与其它 impl 变体的差异、切换 feature 时的影响
% - 性能/内存假设与已知限制
% 【事实来源】
% - 上述 crate 源码与 Cargo.toml
% - os/feature-tree.txt 中指向本 impl 的 feature 链

\subsubsection{impl-kernel}

\paragraph{全局状态。}
\begin{itemize}
 \item \code{ROOT\_FS}：当前根RO\code{SharedFs}，辅助路径；
 \item \code{ROOT\_RW\_FS}：bring-up主路径的RW根卷句柄；
 \item \code{ROOT\_DEV\_PATH}：根块设备路径字符串；
 \item \code{ACTIVE\_FS\_IMPL}：\code{fs::init()}注入的\code{'static dyn FsImpl}；
 \item \code{MOUNT\_GENERATION}：原子计数，每次成功挂载或辅助卷变更后递增，供VFS页缓存失效。
\end{itemize}

\paragraph{挂载流程。}
\code{mount\_root\_rw\_from\_block\_path}：\code{devfs::lookup\_block\_device} 经 \code{ACTIVE\_FS\_IMPL.mount\_rw} 至写入 \code{ROOT\_RW\_FS} 与 \code{ROOT\_DEV\_PATH}，再至 \code{mount\_generation}++。\code{mount\_default\_root\_rw}使用\code{default\_root\_block\_path()}，通常 \code{/dev/vda}。

\paragraph{辅助卷。}
\code{mount\_aux\_ro\_from\_block\_path}/\code{mount\_aux\_rw\_from\_block\_path}创建独立\code{SharedFs}/\code{SharedRwFs}。若目标块设备与当前根卷\code{Arc}相同，RO路径可复用根句柄；RW根已挂载时拒绝同设备辅助RW别名，避免双写歧义。每次辅助挂载调用\code{bump\_mount\_generation}。

\paragraph{对外符号。}
\code{set\_active\_fs\_impl}、\code{root\_rw\_fs}、\code{mount\_generation}、\code{bump\_mount\_generation}由聚合层与VFS bridge直接引用；页缓存键含\code{mount\_gen}，根卷remount或辅助卷变更后须见到代次递增。
